% ЗӨВХӨН БҮТЭН НООРОГТ (main.tex) хэвлэгдэнэ. Багшид үзүүлэх ноорогт (main_teacher.tex) орохгүй.
% 2.2-ын доторх хоёр догол мөр: RAPTOR, Han нар, TaxoIndex, T-RAG (2026.11.16 (4): хэрэглэгч — одоохондоо дараагийн шат;
% хэрэглэгч эдгээрийг уншсаны дараа бүлгийн файл руу буцааж зөөнө).
Модонд суурилсан хайлтын түгээмэл аргын нэг нь RAPTOR юм \cite{sarthi2024raptor}. Энэ арга бичвэрийн хэсгүүдийг автоматаар бүлэглээд, бүлэг бүрийн хураангуйг хэлний загвараар бичиж, модыг доороос дээш босгодог. Хураангуйг хайлтын нэмэлт зангилаа болгон ашигладаг. Энэ ажилд ийм автомат мод босгоогүй нь хоёр шалтгаантай. Нэгдүгээрт, ресепшнд асуудаг сэдвүүд (\tref{tab:reception_topics}) төлбөр, өрөө, холбоо барих хүн зэрэг нэг баримтаар хариулагдах асуулт юм. Ийм асуултад энгийн хайлт дутаагүй: Han нарын үнэлгээнд нэг баримтын асуулттай Natural Questions өгөгдөл дээр энгийн RAG F1 64.78, RAPTOR-т суурилсан арга 60.04 авсан \cite{han2025ragvsgraphrag}. Харин олон баримтыг холбох асуултад модтой арга ялимгүй илүү байсан бөгөөд зохиогчид хоёр төрлийн аргыг бие биеэ нөхдөг гэж дүгнэсэн. Хоёрдугаарт, хэлний загварын бичсэн хураангуй буруу мэдээлэл агуулж болдог. RAPTOR-ын зохиогчид гараар шалгасан 150 зангилааны 4\%-д буруу мэдээлэл байсан \cite{sarthi2024raptor}.

Хүний бэлдсэн ангиллыг хайлтад ашигласан ажил ч бий. TaxoIndex эрдэм шинжилгээний өгүүллийн хайлтад бэлэн салбарын ангиллыг индекст ашигласан \cite{kang2024taxoindex}. T-RAG нь нэг байгууллагын бүтцийн модыг асуултад нэмэлт мэдээлэл болгон ашиглаж, 37 асуултын 27-д зөв хариулсан бол ердийн RAG 21-д зөв хариулсан \cite{fatehkia2024trag}. Эдгээрийн үнэлгээний багц жижиг эсвэл өөр салбарынх тул үр дүнг энэ ажилд шууд шилжүүлэх боломжгүй.
